\clearpage
\SourceLeaf{407}\pdfbookmark[1]{III. Zweite Methode. Mai–Juli 1805}{ceres-second-method}
\begin{center}
{\large BRIEFWECHSEL UND NACHLASS.\par}
\bigskip
[III. ZWEITE METHODE.]\par
[MAI—JULI 1805.]\par
\bigskip
[1.]\par
\medskip
\textsc{Gauss} an \textsc{Olbers}, Braunschweig, 10. Mai 1805.
\end{center}
\noindent\hspace*{\parindent}.\quad.\quad.\quad.\quad.\quad Die Methode, nach der ich die {\AstroSymbols ⚳}-Störungen zu berechnen angefangen hatte, habe ich doch wieder aufgegeben. Das gar zu viele mechanische todte Rechnen, was ich dabei vor mir sah, hat mich abgeschreckt; auch selbst wenn alle Rechnungen, die ich Fremden hätte übertragen können, von Hrn. \textsc{Bessel} und Hrn. v. \textsc{Lindenau}\quad.\quad.\quad.\quad.\quad übernommen wären, würde für mich noch mehr übrig geblieben sein, als meine Geduld hätte bestreiten können.

Ich habe indessen bereits eine andere Methode ausgesonnen, die eben so weit führen kann, als jene, aber bei weitem weniger — obwohl künstlichere — Arbeit erfordert. Ich habe schon stark angefangen, sie auf die Ceres anzuwenden, wiewohl vorerst nur nach einem eingeschränktern Plane, indem ich nur bis an die 5. Potenzen der Excentricitäten von {\AstroSymbols ♃} und {\AstroSymbols ⚳} gehe. Diese Methode hat um so mehr Reiz für mich, da ich dabei von vielen, schon vor längerer Zeit angestellten, ziemlich tiefen Untersuchungen über eigene Arten von transcendenten Functionen einen glücklichen Gebrauch machen kann. Ich werde in [der] Folge Ihnen eine Idee davon zu geben suchen. Auch hoffe ich, dass ich im Stande sein werde, alles so einzurichten, dass ich durch eine fremde Unterstützung eine ansehnliche Erleichterung erhalte, zwar nicht bei meiner diesmaligen Rechnung für die {\AstroSymbols ⚳} (denn gerade in dem Theile, wo Hülfe zu brauchen wäre, bin ich schon selbst zu weit \SourceLeaf{408}vorgerückt), aber doch wenn ich dieselbe wiederhole, welches nöthig sein wird, da ohne Zweifel die erweiterten Störungsgleichungen noch mit ansehnlichen Änderungen in den Elementen selbst verbunden sein werden — oder auch wenn ich einst diese Arbeit bei der {\AstroSymbols ⚴} und {\AstroSymbols ⚵} vornehme, wo sie beträchtlich weitläufiger sein wird.\quad.\quad.\quad.\quad.\quad.
\par\medskip
\begin{center}\rule{.30\linewidth}{.4pt}\end{center}
